\documentclass[11pt]{article}
\usepackage[review]{acl}
\usepackage{times}
\usepackage{latexsym}
\usepackage[T1]{fontenc}
\usepackage[utf8]{inputenc}
\usepackage{microtype}
\usepackage{inconsolata}
\usepackage{graphicx}
\usepackage{booktabs}
\usepackage{amsmath}

\title{TOME: Temporal Opponent Modeling and Outcome-Conditioned Episodic Memory\\
for Adaptive Tactical Agents in Real-Time Strategy Games}
\author{Anonymous ACL submission}

\begin{document}
\maketitle

\begin{abstract}
Real-time strategy (RTS) agents must react to opponents whose tactics change over
a match while acting under strict latency and action-validity constraints.
Existing policies often condition on the current state or optimize against fixed
opponents, making early response to tactical shifts difficult. We propose TOME,
a hierarchical agent that models an opponent from a temporal sequence of game
events, retrieves outcome-conditioned episodic memories from similar prior
games, and uses a low-frequency tactical planner with a constraint-aware action
executor. TOME separates language-model or rule-based deliberation from
tick-level control: the planner selects a discrete tactic and the executor maps
it to permitted actions. This paper presents the method, a pre-registered
MicroRTS evaluation protocol, and explicitly marked pilot measurements. The
pilot results establish the experimental pipeline but do not support a final
performance claim.
\end{abstract}

\section{Introduction}

RTS games require agents to allocate resources, build units, and fight an
adaptive adversary in real time. Large-scale learning has mastered complex
strategic domains \citep{vinyals2019grandmaster,berner2019dota}. However, a
policy trained against a limited opponent distribution can respond poorly when
an opponent changes from economic expansion to an early rush. The useful
evidence is not a single board state, but a sequence of unit production,
movement trends, resource choices, and distances between forces.

Language-model agents make high-level reasoning and tool use practical
\citep{yao2023react,wang2023voyager}, while memory can help agents reuse
experience across tasks \citep{park2023generative,shinn2023reflexion}. Directly
asking an API model to issue unit-level actions is unsuitable for RTS games: it
would be called too frequently, can return invalid actions, and makes latency
and cost part of the control loop. A usable game agent needs a clear boundary
between deliberation and real-time execution.

We study tactical adaptation in MicroRTS \citep{ontanon2013microrts}, using
the \texttt{gym-microrts} interface \citep{huang2021gymmicrorts}. We ask whether
an agent can predict near-future tactical pressure from behavior time series,
retrieve useful prior episodes, and translate a selected tactic into legal
actions quickly enough to affect a match.

\section{Problem Formulation}

At tick $t$, the environment provides observation $o_t$ and legal action mask
$m_t$. TOME controls one player against an opponent with latent tactical mode
$z_t \in \{\textit{economic}, \textit{balanced}, \textit{advance},
\textit{rush}\}$. Every $K$ ticks, it selects a high-level tactic $a^H_t$ from
$\{\textit{economy},\textit{scout},\textit{defend},\textit{rush}\}$. At every
tick, its executor emits a legal low-level action $a^L_t \in \mathcal{A}(m_t)$.
The objective is
\[
\max_{\pi} \; \mathbb{E}[R] - \lambda\mathbb{E}[D_{\mathrm{switch}}],
\]
where $D_{\mathrm{switch}}$ is the delay between an opponent tactical change
and a compatible response. The central challenges are temporal prediction,
selective reuse of delayed-outcome experience, and hard action constraints.

\section{TOME}

\subsection{Temporal Opponent Dynamics Model (TODM)}

TODM converts the recent $W$ observations into an event sequence
$e_{t-W+1:t}$. An event includes own and enemy unit counts by type, worker and
combat-unit growth, resource and building changes, centroid distance, and
movement trend. A temporal encoder predicts a future mode and attack pressure:
\[
(\hat z_{t+H},\hat p_{t+H})=f_\theta(e_{t-W+1:t}).
\]
The current implementation is an interpretable temporal rule model. The final
evaluation will compare it with a single-frame classifier, window statistics,
GRU, TCN, and Transformer models. In particular, worker growth together with
approach toward the agent base is treated as a possible WorkerRush signal,
rather than economic behavior alone.

\subsection{Outcome-Conditioned Episodic Memory (OCEM)}

After episode $j$, OCEM stores
\[
q_j=(\phi_j,\hat z_j,a^H_j,r_j,y_j),
\]
where $\phi_j$ is a state signature, $r_j$ is return, and $y_j \in\{-1,0,1\}$
is the terminal result. At planning time, it retrieves episodes that match the
state signature and predicted opponent mode, ranking them by similarity and
positive outcome. The retrieved tactic is evidence, not an unconditional
command; failed episodes are not treated as positive demonstrations.

\subsection{Constraint-aware Tactical Planner and Executor (CTPE)}

The planner receives a compact structured state summary, a TODM prediction, and
memory evidence. It outputs one token from the tactic vocabulary. It can be a
deterministic rule policy or an API language model, but runs only every $K$
ticks. CTPE runs every tick: it maps a tactic to unit priorities, queries the
MicroRTS action mask, and selects only permitted action types and parameters.
Thus, an API model never directly controls the game loop.

\begin{figure}[t]
\centering
\fbox{\begin{minipage}{0.90\columnwidth}\centering
State sequence $\rightarrow$ TODM $\rightarrow$ OCEM $\rightarrow$ planner
$\rightarrow$ CTPE legal actions $\rightarrow$ outcome $\rightarrow$ memory
\end{minipage}}
\caption{TOME decision loop. The planner runs every $K$ ticks; CTPE runs at
every tick using the action mask.}
\label{fig:tome}
\end{figure}

\section{Research Questions and Expected Contributions}

\paragraph{RQ1: Temporal prediction.} Does temporal opponent modeling improve
early tactical-shift detection and downstream performance over single-frame
state conditioning? We will measure future-mode macro-F1, pressure MAE, and
response delay.

\paragraph{RQ2: Outcome-conditioned memory.} Does retrieving successful
episodes from similar states improve adaptation to unseen opponent mixtures or
strategy switches? We will measure win rate, normalized return, tactic-switch
success, and sample efficiency against no-memory and recency-memory controls.

\paragraph{RQ3: Hierarchical constrained execution.} Does separating
low-rate planning from mask-constrained control preserve real-time feasibility
while improving adaptation? We will measure action validity, latency, API calls,
cost, and performance against a per-tick planning control.

This paper makes three intended contributions: (1) TOME, an RTS agent combining
temporal opponent prediction, outcome-conditioned memory, and constrained
hierarchical execution; (2) an evaluation protocol that jointly measures
prediction, adaptation, action validity, and runtime cost; and (3) systematic
ablations that isolate TODM, OCEM, CTPE, and language-model planning.

\section{Experimental Design}

\subsection{Environment and Opponents}

We will use standard 10$\times$10 MicroRTS base-and-worker maps. The opponent
suite will include RandomAI, WorkerRush, LightRush, MCTS, Minimax/ABCD where
available, single-agent RL, and multi-agent RL policies. We will evaluate both
stationary opponents and scripted within-episode tactic switches. The final
paper will report exact versions, maps, seeds, training budgets, and hardware.

\subsection{Baselines and Ablations}

All methods receive identical observations and action masks. Comparisons are a
fixed tactic, a state-only planner, TOME without TODM, TOME with ordinary
recency memory, TOME without CTPE, and full TOME. Rule and API planners use the
same tactic vocabulary and decision interval. We will run at least three
independent runs with matched seeds and confidence intervals.

\begin{table*}[t]
\centering
\small
\caption{Updated RTS evaluation on the two-worker 10x10 map: three episodes per
opponent (18 games per method), with a 500-tick horizon and 25-tick decision
interval. These are diagnostic results from the optimized unit-aware executor,
not a powered final comparison.}
\label{tab:main}
\begin{tabular}{lccccccc}
\toprule
Method & PassiveAI & RandomAI & WorkerRush & LightRush & Coac & NaiveMCTS & Overall \\
\midrule
Fixed economy & 0/3 & 0/3 & 0/3 & 0/3 & 0/3 & 0/3 & 0/18 \\
Fixed defend & 0/3 & 0/3 & 0/3 & 0/3 & 0/3 & 0/3 & 0/18 \\
Fixed rush & 3/3 & 2/3 & 0/3 & 0/3 & 0/3 & 0/3 & 5/18 \\
TOME w/o TODM & 3/3 & 1/3 & 0/3 & 0/3 & 0/3 & 0/3 & 4/18 \\
TOME (rule planner) & 3/3 & 1/3 & 0/3 & 0/3 & 2/3 & 0/3 & 6/18 \\
TOME (learned TODM) & 3/3 & 3/3 & 0/3 & 0/3 & 0/3 & 0/3 & 6/18 \\
\bottomrule
\end{tabular}
\end{table*}

\begin{table}[t]
\centering
\small
\caption{Initial RQ1 results on 535 held-out trajectory windows. Response delay
is non-negative delay for detecting a labeled rush, in ticks.}
\label{tab:rq1}
\begin{tabular}{lccc}
\toprule
Model & Mode macro-F1 & Pressure MAE & Response delay \\
\midrule
Single frame & 0.377 & 0.049 & 22.2 \\
Window statistics & 0.515 & 0.050 & 26.7 \\
Temporal sequence (linear) & 0.358 & 0.047 & 111.1 \\
\bottomrule
\end{tabular}
\end{table}

\begin{table}[t]
\centering
\small
\caption{API feasibility pilot on one 100-tick RandomAI episode per row.
Provider pricing was not returned, so cost is not reported.}
\label{tab:efficiency}
\begin{tabular}{lccc}
\toprule
Planner & Latency (ms) & API calls/episode & Cost/episode \\
\midrule
Rule & 0.01 & 0 & -- \\
API, $K=25$ & 1080.98 & 4 & -- \\
API, $K=50$ & 1257.50 & 2 & -- \\
API, $K=100$ & 1230.39 & 1 & -- \\
\bottomrule
\end{tabular}
\end{table}

\section{Pilot Results and Current Implementation Status}

Table~\ref{tab:main} verifies that the optimized, mask-constrained CTPE has zero
invalid requests in the current 90-game diagnostic suite. TOME's rule planner
wins 6/18 games (33.3\%), compared with 4/18 (22.2\%) without TODM and 5/18
(27.8\%) for a fixed rush policy. The gain is concentrated on PassiveAI and
Coac; all methods still fail against WorkerRush and LightRush. This remains a
prediction-to-control gap, not evidence that the current temporal predictor is
solved.
This negative result identifies the hand-written prediction-to-tactic mapping,
rather than the action-validity layer, as the immediate research bottleneck.
In Table~\ref{tab:rq1}, window statistics had the strongest held-out macro-F1
(0.515), but no temporal baseline reduced rush detection delay. This motivates
learned temporal encoders and controlled switch data rather than a claim that
sequence modeling is already solved.

The current prototype implements a game-event decoder, a rule-based TODM,
persistent outcome-conditioned memory, a four-tactic planning interface, and a
unit-aware mask-constrained executor that maps Base, Barracks, Worker, and
combat-unit actions separately. It has been evaluated on all six built-in
opponents using the two-worker 10x10 map. The API feasibility condition makes
one high-level request per decision interval, with CTPE maintaining legal
tick-level control.
The current learned models are NumPy temporal feature baselines. Deep temporal
encoders, RL and MARL opponents, controlled within-episode opponent switches,
and a properly powered final evaluation remain pending.

\section{Limitations and ACL Relevance}

MicroRTS is a simplified RTS environment, so transfer to commercial games is
not guaranteed. Outcome-conditioned memory may amplify biased early experience.
The API condition requires model version, prompt, timeout, failure rate, cost,
and reproducibility reporting. For ACL relevance, language is treated as a
substantive experimental variable: structured event serialization, retrieved
evidence, and constrained tactic generation are explicit planner inputs and
outputs. This comparison should remain in the final paper only if it yields a
clear language-related finding beyond the game-engineering contribution.

\section{Conclusion}

TOME frames RTS tactical adaptation as temporal opponent modeling, selective
reuse of outcome-conditioned experience, and constrained hierarchical control.
The current draft fixes the problem statement, method, research questions, and
evaluation protocol, and reports a reproducible diagnostic suite. The next
milestone is a powered multi-seed evaluation plus learned temporal encoders and
explicit within-episode opponent switches.

\bibliographystyle{acl_natbib}
\bibliography{custom,anthology}
\end{document}
